\section{Integrated Dipoles and Virtual Cancellation}

The subtraction term $d\sigma^A$ is constructed to reproduce the singular
behavior of the real-emission contribution. In order to combine it with
the virtual corrections, it must be integrated analytically over the
unresolved phase space.

\subsection{Integrated subtraction term}

The integration of the subtraction term over the unresolved emission
leads to
\begin{align}
\int_1 d\sigma^A = d\sigma^B \otimes \mathbf{I}.
\end{align}

Here, $\mathbf{I}$ is an operator acting on the Born matrix element,
and it contains explicit poles in dimensional regularization:
\begin{align}
\mathbf{I} \sim \frac{1}{\epsilon^2} + \frac{1}{\epsilon}.
\end{align}

These poles match exactly the infrared divergences of the virtual
contribution $d\sigma^V$.

\subsection{Cancellation of divergences}

The virtual contribution has the structure
\begin{align}
d\sigma^V \sim \frac{1}{\epsilon^2} + \frac{1}{\epsilon} + \text{finite}.
\end{align}

When combined with the integrated subtraction term, the poles cancel:
\begin{align}
d\sigma^V + \int_1 d\sigma^A = \text{finite}.
\end{align}

This cancellation ensures that the virtual contribution can be evaluated
in four dimensions.

\subsection{Initial-state collinear contributions}

For hadronic processes, additional collinear singularities arise from
initial-state radiation. These are absorbed into the parton distribution
functions (PDFs) through mass factorization.

This introduces additional terms, commonly denoted as $\mathbf{P}$ and
$\mathbf{K}$ operators:
\begin{align}
\sigma^{\text{NLO}} \supset \int dx \, \left( \mathbf{P} + \mathbf{K} \right).
\end{align}

\begin{itemize}
    \item The $\mathbf{P}$ term is related to Altarelli--Parisi splitting
    functions and governs the scale dependence of PDFs.
    \item The $\mathbf{K}$ term contains finite remainders from the
    factorization procedure.
\end{itemize}

\subsection{Interpretation}

The integrated dipoles complete the subtraction procedure:
\begin{itemize}
    \item they cancel the infrared poles of the virtual corrections,
    \item they account for initial-state collinear singularities,
    \item they ensure that the full NLO result is finite.
\end{itemize}